\chapter{Young 方法的结构定理与标准基}
\label{app:young-proofs}

\chapref{chap:permutation-group}已经建立 Young 盘、行列群与 Specht 模的基本构造。本章接着证明这些构造为何给出全部不可约表示，并把标准基与钩长公式的计算依据逐一补齐。每个结论先明确假设，再给出独立证明；这些结构定理也为下一章的特征标计算提供输入。沿用该章记号
\[
a_T=\sum_{r\in R_T}r,
\qquad
b_T=\sum_{c\in C_T}\operatorname{sgn}(c)c,
\qquad
 e_T=a_Tb_T.
\]

\section{幂等元、左理想与不可约性}

\begin{theorem}[本原幂等元判据]
设 $A=\mathbb C[G]$ 为有限群的复群代数，$p\in A$ 为非零幂等元 $p^2=p$。左理想 $Ap$ 是不可约左 $A$-模，当且仅当
\begin{equation}
pAp=\CC p,
\label{eq:primitive-idempotent-criterion}
\end{equation}
即对任意 $x\in A$ 都存在 $\lambda_x\in\CC$ 使
\[
pxp=\lambda_xp.
\]

\end{theorem}

\begin{proof}
先设 $Ap$ 不可约。对任意 $x\in A$，定义
\[
T_x:Ap\to Ap,
\qquad
T_x(ap)=apxp.
\]
它确实把 $Ap$ 映到自身，并且对任意 $y\in A$，
\[
T_x(y\cdot ap)=yapxp=y\cdot T_x(ap),
\]
所以 $T_x$ 是 $A$-模自同态。由 Schur 引理，存在 $\lambda_x$ 使
\[
T_x=\lambda_x I_{Ap}.
\]
特别取向量 $p\in Ap$，得到
\[
pxp=\lambda_xp,
\]
故 $pAp=\CC p$。

反过来设 \eqnrefs{eq:primitive-idempotent-criterion} 成立。若 $Ap$ 可约，由有限群代数的半单性，存在非零真子模 $W\subsetneq Ap$ 及其不变互补，因此存在与左乘作用交换的非平凡投影
\[
P:Ap\to Ap,
\qquad
P^2=P,
\qquad P\neq0,I.
\]
令 $q=P(p)\in Ap$。因为 $P$ 是 $A$-线性的，对任意 $a\in A$，
\[
P(ap)=aP(p)=aq.
\]
又因 $p^2=p$，
\[
q=P(p)=P(p^2)=pP(p)=pq,
\]
而 $q\in Ap$ 可写为 $q=xp$，所以
\[
q=pxp\in pAp.
\]
按假设 $q=\lambda p$。于是
\[
P(ap)=a\lambda p=\lambda ap,
\]
即 $P=\lambda I$。投影条件给出 $\lambda^2=\lambda$，故 $\lambda=0$ 或 $1$，与 $P$ 非平凡矛盾。因此 $Ap$ 不可约。
\end{proof}

这个判据把“不可约”转化为一个可以在群代数中直接验证的乘法条件。Young 理论后面真正要证明的，正是适当归一化后的 $e_T$ 满足这个条件。

\section{同形 Young 盘之间的共轭关系}

\begin{lemma}[行群与列群随重标号共轭]
若两个 Young 盘满足 $T'=\pi T$，其中 $\pi\in S_n$ 逐格重标号，则
\begin{align}
R_{T'}&=\pi R_T\pi^{-1},\\
C_{T'}&=\pi C_T\pi^{-1},\\
a_{T'}&=\pi a_T\pi^{-1},\\
b_{T'}&=\pi b_T\pi^{-1},\\
e_{T'}&=\pi e_T\pi^{-1}.
\label{eq:young-conjugation-all}
\end{align}

\end{lemma}

\begin{proof}
若 $r\in R_T$，它只交换 $T$ 的同一行中的数字。重标号后，$\pi r\pi^{-1}$ 对 $T'$ 做完全相同的格位置交换，因此属于 $R_{T'}$。反向使用 $\pi^{-1}$ 得到集合相等。列群同理。

对群代数求和，
\[
\pi a_T\pi^{-1}
=\sum_{r\in R_T}\pi r\pi^{-1}
=\sum_{r'\in R_{T'}}r'=a_{T'}.
\]
共轭不改变置换奇偶性，因此 $b_T$ 也同样变换。最后两式相乘即得 $e_{T'}=\pi e_T\pi^{-1}$。
\end{proof}

\begin{proposition}[同形盘给出的左理想同构]
若 $T'=\pi T$，则右乘 $\pi^{-1}$ 给出左 $A$-模同构
\[
Ae_T\longrightarrow Ae_{T'},\qquad xe_T\longmapsto xe_T\pi^{-1}.
\]
\end{proposition}

\begin{proof}
同形不同盘对应的左理想不会产生新的不可约等价类。由
\[
e_{T'}=\pi e_T\pi^{-1}
\]
先得到
\[
Ae_{T'}=A\pi e_T\pi^{-1}=Ae_T\pi^{-1},
\]
其中使用了 $A\pi=A$。因此可以定义
\[
\Phi_\pi:Ae_T\to Ae_{T'},
\qquad
\Phi_\pi(xe_T)=xe_T\pi^{-1}=x\pi^{-1}e_{T'}.
\]
若 $xe_T=x'e_T$，右乘 $\pi^{-1}$ 后仍相等，所以映射良定义。对任意 $a\in A$，
\[
\Phi_\pi(a\,xe_T)=axe_T\pi^{-1}=a\,\Phi_\pi(xe_T),
\]
故它与左正则作用交织。右乘 $\pi$ 给出逆映射，于是 $Ae_T\simeq Ae_{T'}$。因此“表示由图形状而不是数字填法标记”是一个明确的模同构结论。
\end{proof}

\section{行列冲突与不同形状的正交性}

Young 算符的零乘积并不是偶然的系数抵消，而是“同行要求对称、同列要求反对称”之间的直接冲突。先把这个局部机制写成一个可反复调用的引理。

\begin{lemma}[行列冲突引理]
设 $T,S$ 为两个 Young 盘。若存在两个数字 $a,b$ 在 $T$ 中位于同一行，而在 $S$ 中位于同一列，则
\begin{equation}
b_Sa_T=0.
\label{eq:row-column-collision}
\end{equation}
从而
\[
e_Se_T=a_Sb_Sa_Tb_T=0.
\]

\end{lemma}

\begin{proof}
令 $\tau=(ab)$。因为 $a,b$ 在 $T$ 同行，所以 $\tau\in R_T$，由行群重排有
\[
\tau a_T=a_T.
\]
因为它们在 $S$ 同列，所以 $\tau\in C_S$，且 $\operatorname{sgn}(\tau)=-1$。列反对称子的定义给出
\[
b_S\tau=-b_S.
\]
于是
\[
b_Sa_T=b_S(\tau a_T)=(b_S\tau)a_T=-b_Sa_T,
\]
故 $2b_Sa_T=0$。我们工作在复数域上，特征不为 $2$，因此只能有 $b_Sa_T=0$。
\end{proof}

同形盘还需要一个互补结论：若没有行列冲突，那么相应的重标号恰能分解成“行置换乘列置换”。这一步正是后面证明 $e_Txe_T$ 一维性的关键。

\begin{lemma}[同形盘的行列分解判据]
设 $T$ 与 $T'=\pi T$ 具有相同形状。若 $T$ 中同一行的任意两个数字，在 $T'$ 中都不位于同一列，则存在
\[
r\in R_T,\qquad c\in C_T
\]
使
\begin{equation}
\pi=rc.
\label{eq:row-column-factorization}
\end{equation}
反过来，若 $\pi=rc$，则不会出现上述行列冲突。因此
\[
\pi\notin R_TC_T
\quad\Longleftrightarrow\quad
\text{$T$ 某两个同行数字在 $\pi T$ 中同列}.
\]

\end{lemma}

\begin{proof}
先假设不存在行列冲突。考察 $T$ 的第一行。它含 $\lambda_1$ 个数字，而形状 $\lambda$ 一共恰有 $\lambda_1$ 列；这些数字在 $T'$ 中落在互不相同的列，所以恰好每列一个。因而可以只在 $T$ 第一行内部重排数字，使第一行每个数字都进入它在 $T'$ 中所在的列。

删除已经配好的第一行数字后，仍含未配数字的列数正好是 $\lambda_2$。$T$ 第二行含 $\lambda_2$ 个数字；无行列冲突保证它们在 $T'$ 的剩余位置中仍落在互异列，因此又可以只在第二行内部重排，使这些数字进入正确的列。逐行重复，得到一个行置换 $r\in R_T$，使 $rT$ 与 $T'$ 的每一列包含完全相同的一组数字，只是列内上下顺序可能不同。

于是存在 $rT$ 的列置换 $c'\in C_{rT}$ 使
\[
T'=c'rT.
\]
由重标号共轭关系 $C_{rT}=rC_Tr^{-1}$，可写 $c'=rcr^{-1}$，其中 $c\in C_T$。因此
\[
T'=rcr^{-1}rT=rcT.
\]
Young 盘中每个数字只出现一次，所以若两个置换作用在 $T$ 上得到同一盘，它们必相等，故 $\pi=rc$。

反过来设 $\pi=rc$。若出现两个数字在 $T$ 同行而在 $\pi T$ 同列，令交换这两个数字的对换为 $\tau$，则
\[
\tau\in R_T\cap \pi C_T\pi^{-1}.
\]
由 $\pi=rc$ 以及 $cC_Tc^{-1}=C_T$，有
\[
\pi C_T\pi^{-1}=rC_Tr^{-1}.
\]
于是 $r^{-1}\tau r\in R_T$（因为 $r\in R_T$ 且 $R_T$ 是子群），同时 $r^{-1}\tau r\in C_T$，所以
\[
r^{-1}\tau r\in R_T\cap C_T=\{e\}.
\]
这迫使 $\tau=e$，与 $\tau$ 是交换两个不同数字的对换矛盾。因此不存在行列冲突。
\end{proof}

现在比较不同形状。这里使用\textbf{字典序}：若
\[
\lambda=(\lambda_1,\lambda_2,\dots),\qquad
\mu=(\mu_1,\mu_2,\dots)
\]
在第一个不同的行指标 $k$ 满足 $\lambda_k>\mu_k$，就记作 $\lambda\succ\mu$。任意两个不同分拆在这个顺序下都可比较。

\begin{lemma}[不同形状必产生行列冲突]
若 $\lambda\succ\mu$，$T$ 是形状 $\lambda$ 的任意 Young 盘，$S$ 是形状 $\mu$ 的任意 Young 盘，则必存在两个数字在 $T$ 中同行而在 $S$ 中同列。

\end{lemma}

\begin{proof}
反设不存在这样的数字。由于 $T$ 第一行的 $\lambda_1$ 个数字在 $S$ 中必须占据互不相同的列，而 $S$ 一共只有 $\mu_1$ 列，所以
\[
\lambda_1\le\mu_1.
\]
若第一行已经是两个分拆首次不同之处，这就与 $\lambda_1>\mu_1$ 矛盾。

若 $\lambda_1=\mu_1$，那么 $T$ 第一行的数字必须在 $S$ 的每一列中恰出现一个。把这些数字暂时删去。此时仍含未删数字的列恰有 $\mu_2$ 个，因为原来列高至少为 $2$ 的列数就是第二行长度 $\mu_2$。$T$ 第二行的 $\lambda_2$ 个数字仍必须落在互不相同的剩余列中，所以 $\lambda_2\le\mu_2$。

若前两行仍相等，就继续删除已经匹配的数字。归纳下去：只要
\[
\lambda_1=\mu_1,\dots,\lambda_{j-1}=\mu_{j-1},
\]
删除前 $j-1$ 行对应数字以后，$S$ 中仍有未删数字的列数恰为 $\mu_j$，而 $T$ 第 $j$ 行的 $\lambda_j$ 个数字必须占据不同列，因此 $\lambda_j\le\mu_j$。取第一个不同指标 $k$ 就得到 $\lambda_k\le\mu_k$，与 $\lambda_k>\mu_k$ 矛盾。
\end{proof}

\begin{proposition}[异形 Young 算符的一向零乘积]
若 $T,S$ 分别具有形状 $\lambda,\mu$，且 $\lambda\succ\mu$，则
\[
e_Sxe_T=0,\qquad \forall x\in A=\CC[S_n].
\]
\end{proposition}

\begin{proof}
组合引理立即给出群代数中的零乘积。设 $\lambda\succ\mu$，$T,S$ 分别具有形状 $\lambda,\mu$。对任意 $\pi\in S_n$，盘 $\pi T$ 仍具有形状 $\lambda$，所以由上个引理，$\pi T$ 的某一行与 $S$ 的某一列发生冲突，从而
\[
b_Sa_{\pi T}=0.
\]
另一方面
\[
a_{\pi T}=\pi a_T\pi^{-1},
\qquad
\pi a_T=a_{\pi T}\pi.
\]
因此
\begin{align}
e_S\pi e_T
&=a_Sb_S\pi a_Tb_T\notag\\
&=a_Sb_Sa_{\pi T}\pi b_T=0.
\label{eq:different-shape-zero}
\end{align}
线性延拓后，对任意 $x\in\mathbb C[S_n]$ 都有
\[
e_Sxe_T=0.
\]
这条式子只需要在字典序较大形状 $\lambda$ 到较小形状 $\mu$ 的一个方向成立；它已经足以排除两个 Specht 模之间存在同构，因为若存在同构，其逆映射会给出该方向上的非零模同态。
\end{proof}

\section{\texorpdfstring{Young 对称化子的本原性}{Young 对称化子的本原性}}

\begin{lemma}[Young 一维性引理]
固定 Young 盘 $T$。若 $x\in\mathbb C[S_n]$ 满足
\begin{equation}
rxc=\operatorname{sgn}(c)x,
\qquad
\forall r\in R_T,\ c\in C_T,
\label{eq:young-bi-equivariance}
\end{equation}
则 $x$ 必为 $e_T$ 的标量倍数。

\end{lemma}

\begin{proof}
把
\[
x=\sum_{\pi\in S_n}x_\pi\pi
\]
按置换基展开。先证明若 $\pi\notin R_TC_T$，则 $x_\pi=0$。由 \eqnrefs{eq:row-column-factorization}，此时存在两个数字在 $T$ 中同行而在 $\pi T$ 中同列。令交换这两个数字的对换为 $\tau$。于是
\[
\tau\in R_T,
\qquad
\tau\in C_{\pi T}=\pi C_T\pi^{-1}.
\]
因此存在一个对换 $c\in C_T$ 使
\[
\tau=\pi c\pi^{-1},
\qquad
\tau\pi=\pi c.
\]
由左侧行不变性 $\tau x=x$，比较基向量 $\tau\pi$ 的系数得到
\[
x_{\tau\pi}=x_\pi.
\]
由右侧列反对称性 $xc=-x$，比较基向量 $\pi c$ 的系数得到
\[
x_{\pi c}=-x_\pi.
\]
但 $\tau\pi=\pi c$，故 $x_\pi=-x_\pi$，于是 $x_\pi=0$。

所以 $x$ 的支撑只能落在集合 $R_TC_T$ 上。若
\[
rc=r'c',
\qquad r,r'\in R_T,\quad c,c'\in C_T,
\]
则
\[
r^{-1}r'=cc'^{-1}\in R_T\cap C_T=\{e\},
\]
故分解 $rc$ 唯一。设单位元系数 $x_e=\theta$。由 $rx=x$ 可知左乘行置换不改变系数；由 $xc=\operatorname{sgn}(c)x$ 可知右乘列置换只给出相应符号，因此
\[
x_{rc}=\theta\,\operatorname{sgn}(c).
\]
于是
\[
x
=\theta\sum_{r\in R_T}\sum_{c\in C_T}\operatorname{sgn}(c)rc
=\theta a_Tb_T
=\theta e_T.
\]
\end{proof}

\begin{theorem}[Young 准幂等元与本原归一化]
对形状为 $\lambda$ 的 Young 盘 $T$，存在只依赖形状的非零常数 $c_\lambda$，使
\[
e_T^2=c_\lambda e_T,\qquad
p_T=c_\lambda^{-1}e_T,\qquad
p_T^2=p_T,\qquad p_TAp_T=\CC p_T.
\]
因此 $Ap_T=Ae_T$ 是不可约左 $A$-模。并且
\[
c_\lambda=\frac{n!}{\dim Ae_T}.
\]
\end{theorem}

\begin{proof}
先令 $x=e_Tye_T$。由于 $r a_T=a_T$ 以及 $b_Tc=\operatorname{sgn}(c)b_T$，它自动满足 \eqnrefs{eq:young-bi-equivariance}，于是
\begin{equation}
e_Tye_T=\lambda_T(y)e_T,
\qquad \forall y\in\mathbb C[S_n].
\label{eq:ete-scalar}
\end{equation}
特别取 $y=e$，得到
\[
e_T^2=c_\lambda e_T.
\]
这给出了\chapref{chap:permutation-group}所谓“本质幂等性”，但归一化以前必须证明 $c_\lambda\ne0$。
在 $e_T=a_Tb_T$ 的置换基展开中，单位元只由 $r=c=e$ 产生，
因为 $R_T\cap C_T=\{e\}$，所以它的系数为 $1$。
右乘算符 $R_{e_T}:A\to A$ 的每个置换基向量都具有相同的对角系数，故
\[
\operatorname{tr}R_{e_T}=n!,\qquad R_{e_T}^2=c_\lambda R_{e_T}.
\]
若 $c_\lambda=0$，则 $R_{e_T}$ 为平方零算符，迹必为 $0$，矛盾。
同形盘共轭又保证这个常数只依赖形状。现在才可以归一化
\[
p_T=c_\lambda^{-1}e_T
\]
以后，\eqnrefs{eq:ete-scalar} 变成
\[
p_Typ_T=\mu_T(y)p_T.
\]
由本原幂等元判据，$\mathbb C[S_n]p_T$ 不可约。
此外，$R_{p_T}$ 是像空间为 $Ap_T$ 的幂等投影，所以
\[
\dim Ae_T=\operatorname{tr}R_{p_T}=\frac{n!}{c_\lambda}.
\]
这同时得到常数的表达式。
\end{proof}

\section{同形等价、异形不等价与完备性}

\begin{theorem}[Specht 模的互异性与完备性]
所有 $\lambda\vdash n$ 对应的 $Ae_T$ 构成 $S_n$ 的全部复不可约表示，
其同构类只依赖 $T$ 的形状；不同形状给出不同同构类。
\end{theorem}

\begin{proof}
同形等价已经由 \eqnrefs{eq:young-conjugation-all} 及右乘 $\pi^{-1}$ 的交织映射得到。现在设 $T,S$ 形状分别为 $\lambda,\mu$，且 $\lambda\succ\mu$。前节已经证明
\[
e_Sxe_T=0,
\qquad \forall x\in A.
\]
将本质幂等元归一化为
\[
p_T=c_\lambda^{-1}e_T,
\qquad
p_S=c_\mu^{-1}e_S,
\]
则 $Ap_T=Ae_T$、$Ap_S=Ae_S$，并且 $p_T^2=p_T$、$p_S^2=p_S$。若这两个左 $A$-模等价，则存在模同构 $\Phi:Ap_T\to Ap_S$，从而逆映射
\[
\Psi=\Phi^{-1}:Ap_S\to Ap_T
\]
也是非零 $A$-模同态。它由 $\Psi(p_S)$ 决定。利用幂等性与 $A$-线性，
\[
\Psi(p_S)=\Psi(p_S^2)=p_S\Psi(p_S).
\]
另一方面 $\Psi(p_S)\in Ap_T$，故存在某个 $y\in A$ 使
\[
\Psi(p_S)=yp_T.
\]
两式结合得到
\[
\Psi(p_S)=p_Syp_T
=\frac{1}{c_\mu c_\lambda}e_Sye_T=0,
\]
这与 $\Psi$ 为同构矛盾。因此不同形状的 Specht 模不可能等价。

最后，$S_n$ 的共轭类由 $n$ 的整数分拆分类，因此共轭类数正好是分拆数 $p(n)$。有限群复不可约表示等价类数又等于共轭类数。我们已经对每个 $\lambda\vdash n$ 构造出一个不可约且彼此不等价的 Specht 模，所以这些模必然已经穷尽全部不可约表示。这一步把“构造出很多表示”升级成了“分类完毕”。
\end{proof}

\section{标准列反对称盘向量基与 Specht 模维数}

正文用群代数左理想 $Ae_T$ 构造不可约表示；计算维数时，经典行等价盘模型更透明。这里把两种语言之间的关系说明清楚，并证明“标准 Young 盘数就是不可约表示维数”。

\begin{definition}[行等价盘模]
固定分拆 $\lambda\vdash n$。若两个 $\lambda$-盘只相差同行内部的重排，
就把它们视为同一个行等价盘，记为 $\{T\}$。
全部行等价盘构成的复向量空间记为 $M^\lambda$，$S_n$ 通过重标号线性作用。
\end{definition}

固定参考盘 $T$ 后，映射
\[
A a_T\longrightarrow M^\lambda,
\qquad
x a_T\longmapsto x\{T\}
\]
是良定义的 $A$-模同构：右乘任意 $r\in R_T$ 不改变 $a_T$，也不改变行等价盘；反过来，不同左陪集对应不同行等价盘。因此 $M^\lambda$ 正是行群平凡表示诱导出的置换模。

对任意 $\lambda$-盘 $T$ 定义列反对称盘向量
\begin{equation}
\varepsilon_T
=\sum_{c\in C_T}\operatorname{sgn}(c)\{cT\}.
\label{eq:polytabloid-app}
\end{equation}
其张成空间记为 $\mathscr S^\lambda\subset M^\lambda$。它在 $S_n$ 作用下不变，因为
\[
\pi\varepsilon_T=\varepsilon_{\pi T}.
\]
这就是经典 Specht 模的行等价盘实现。

\begin{lemma}[Garnir 拉直关系]
设 $T$ 的每列已经严格递增，但第 $j,j+1$ 列在第 $i$ 行出现行逆序。
记这两列的数字为 $a_1<\cdots<a_p$、$b_1<\cdots<b_q$，且 $a_i>b_i$。
取列带
\[
A=\{a_i,\dots,a_p\},\qquad B=\{b_1,\dots,b_i\}.
\]
令 $\mathcal D(A,B)$ 为 $S_A\times S_B$ 在 $S_{A\cup B}$ 中的左陪集代表，
即分解 $S_{A\cup B}=\bigsqcup_{\sigma\in\mathcal D}\sigma(S_A\times S_B)$；
可选取把各组重新填回对应列段时保持组内递增的保序交错代表，并包括恒等元。
则
\begin{equation}
\sum_{\sigma\in\mathcal D(A,B)}\operatorname{sgn}(\sigma)\,
\varepsilon_{\sigma T}=0.
\label{eq:garnir-straightening-app}
\end{equation}
非恒等项经列内排序后，其从左到右、每列从上到下的列读字严格小于 $T$。
\end{lemma}

\begin{proof}
记 $\alpha_E=\sum_{\tau\in S_E}\operatorname{sgn}(\tau)\tau$。
因为 $A,B$ 分别位于 $T$ 的两列，列反对称性给出
\[
\alpha_A\alpha_B\varepsilon_T=|A|!|B|!\varepsilon_T.
\]
陪集分解则给出
\[
\alpha_{A\cup B}
=\Bigl(\sum_{\sigma\in\mathcal D}\operatorname{sgn}(\sigma)\sigma\Bigr)
\alpha_A\alpha_B.
\]
另一方面，$|A\cup B|=p+1$，而在任何 $\{cT\}$、$c\in C_T$ 中，
这些数字都仍处于前 $p$ 行。因此其中必有两个数字同行。
交换它们的对换 $\tau$ 固定这个行等价盘，故
\[
\alpha_{A\cup B}\{cT\}
=\alpha_{A\cup B}\tau\{cT\}
=-\alpha_{A\cup B}\{cT\}=0.
\]
逐项相加，得到 $\alpha_{A\cup B}\varepsilon_T=0$；除以非零整数
$|A|!|B|!$ 就是 \eqnrefs{eq:garnir-straightening-app}。

单调性来自 $b_1<\cdots<b_i<a_i<\cdots<a_p$。
非恒等保序交错置换必把至少一个 $B$ 中较小数字移入第 $j$ 列，
并把相同数目的 $A$ 中较大数字移出；该列排序后的第一个变化位置严格减小。
前面各列不变，所以列读字严格减小。对其余列重新排序只改变
列反对称盘向量的一个符号，且不改变这个比较。
\end{proof}

\begin{proposition}[标准列反对称盘向量的张成性]
\[
\mathscr S^\lambda
=\Span\{\varepsilon_T\mid T\text{ 为标准 }\lambda\text{-盘}\}.
\]
\end{proposition}

\begin{proof}
首先在每列内排序。对 $c\in C_T$ 有
$\varepsilon_{cT}=\operatorname{sgn}(c)\varepsilon_T$，因此这一步只改变符号。
若盘还不标准，它必有相邻两列的某个行逆序，应用 Garnir 关系即可
把它换成列读字更小的盘。盘数有限，严格递降过程必终止；
终点同时行、列递增，只剩标准盘。
注意不能先任意同行排序：行等价盘虽忽略同行次序，列反对称盘向量中的
列群却会随填法改变，$\varepsilon_T$ 不因此保持不变。
\end{proof}

\begin{lemma}[标准列反对称盘向量线性无关]
不同标准盘对应的 $\varepsilon_T$ 线性无关。

\end{lemma}

\begin{proof}
需要把“首项不同”说得足够精确。对任意行等价盘$X$、整数 $k=1,\dots,n$ 和前 $r$ 行，记
\[
N_{k,r}(X)=\#\{1,\dots,k\}\text{ 中位于 }X\text{ 前 }r\text{ 行的数字个数}.
\]
若 $T$ 是标准盘，则在它的每一列中，较小的数字总在较大的数字上方。因此固定 $k,r$ 后，在该列包含的数字中，把 $\{1,\dots,k\}$ 内的那些数字放在尽可能靠上的位置，正是 $T$ 自己所做的排列。任何列置换 $c\in C_T$ 都不能增加前 $r$ 行中小于等于 $k$ 的数字数目，所以
\begin{equation}
N_{k,r}(\{cT\})\le N_{k,r}(\{T\}),
\qquad \forall k,r.
\label{eq:tabloid-dominance-triangular}
\end{equation}
若 $c\neq e$，至少有一列的相对次序被改变，于是对某组 $k,r$ 上式严格小于。

现在在标准行等价盘上定义偏序：若对所有 $k,r$ 都有 $N_{k,r}(S)\le N_{k,r}(T)$，就记 $S\preceq T$。取这个偏序的任一线性扩张作为全序。由 \eqnrefs{eq:polytabloid-app}，$\varepsilon_T$ 中 $c=e$ 的项是 $\{T\}$，系数为 $1$；所有 $c\neq e$ 产生的行等价盘若恰好与某个标准行等价盘行等价，则都严格位于 $T$ 之前。于是把标准列反对称盘向量的坐标投影到标准行等价盘所张成的坐标子空间后，系数矩阵为单位三角矩阵。单位三角矩阵可逆，所以标准列反对称盘向量必线性无关。
\end{proof}

\begin{theorem}[标准基与两种 Specht 实现的维数]
标准列反对称盘向量构成 $\mathscr S^\lambda$ 的一组基，并且
\[
\dim\mathscr S^\lambda=\dim Ae_T=f^\lambda.
\]
\end{theorem}

\begin{proof}
张成性与线性无关性合起来得到经典标准基定理
\begin{equation}
\dim\mathscr S^\lambda=f^\lambda.
\label{eq:standard-polytabloid-dimension-app}
\end{equation}
还需把它与正文使用的 $Ae_T=Aa_Tb_T$ 联系起来。定义群代数自同构
\[
\omega:A\to A,
\qquad
\omega(g)=\operatorname{sgn}(g)g.
\]
若 $T^{\mathsf T}$ 为转置 Young 盘，则 $R_T=C_{T^{\mathsf T}}$、$C_T=R_{T^{\mathsf T}}$，从而
\[
\omega(a_T)=b_{T^{\mathsf T}},
\qquad
\omega(b_T)=a_{T^{\mathsf T}},
\qquad
\omega(e_T)=b_{T^{\mathsf T}}a_{T^{\mathsf T}}.
\]
右端正是经典 Specht 构造的 Young symmetrizer。扭曲一个模的作用 by $\omega$ 不改变向量空间维数，因此正文的 $Ae_T$ 与转置形状的经典 Specht 模具有相同维数。标准盘转置给出标准 $\lambda$-盘与标准 $\lambda^{\mathsf T}$-盘的一一对应，所以
\[
\dim Ae_T=f^{\lambda^{\mathsf T}}=f^\lambda.
\]
这就补全了正文标准盘基定理所需的维数论证。
\end{proof}

\section{Hook-length 公式的递归闭合}

\chapref{chap:permutation-group}已经给出标准盘数的分支递推
\begin{equation}
f^\lambda=\sum_{\mu\nearrow\lambda}f^\mu,
\label{eq:hook-recursion-app}
\end{equation}
其中 $\mu\nearrow\lambda$ 表示从 $\lambda$ 删除一个可移去角格。本节严格证明钩长乘积
\[
H_\lambda=\prod_{u\in\lambda}h_\lambda(u)
\]
满足与 \eqnrefs{eq:hook-recursion-app} 相容的递推，从而闭合正文中的归纳证明。

先引入一个只涉及 Young 图几何的乘积恒等式。设
\[
\lambda=(\lambda_1,\dots,\lambda_r),
\qquad
\lambda_1\ge\cdots\ge\lambda_r>0,
\]
并定义严格递减的整数
\begin{equation}
\ell_i=\lambda_i+r-i,
\qquad i=1,\dots,r.
\label{eq:shifted-row-lengths}
\end{equation}

\begin{lemma}[钩长乘积的移位行长形式]
有
\begin{equation}
H_\lambda
=\frac{\displaystyle\prod_{i=1}^{r}\ell_i!}
{\displaystyle\prod_{1\le i<j\le r}(\ell_i-\ell_j)}.
\label{eq:hook-product-shifted}
\end{equation}

\end{lemma}

\begin{proof}
对行数 $r$ 归纳。$r=1$ 时，Young 图只有一行，hook lengths 为 $\lambda_1,\lambda_1-1,\dots,1$，所以 $H_\lambda=\lambda_1!=\ell_1!$。

设 $r>1$，删去第一行得到
\[
\widehat\lambda=(\lambda_2,\dots,\lambda_r).
\]
删去上方一行不会改变原来第 $2,\dots,r$ 行各格右侧和下方的格数，因此这些格子的 hook lengths 与 $\widehat\lambda$ 中完全相同。于是
\[
H_\lambda=H_{\widehat\lambda}\prod_{j=1}^{\lambda_1}h_{1j}.
\]
对 $\widehat\lambda$ 而言，它的移位行长仍恰为 $\ell_2,\dots,\ell_r$，故归纳假设给出
\[
H_{\widehat\lambda}
=\frac{\prod_{i=2}^{r}\ell_i!}
{\prod_{2\le i<j\le r}(\ell_i-\ell_j)}.
\]

第一行第 $j$ 格的 hook length 为
\[
h_{1j}=\lambda_1-j+\lambda'_j,
\]
其中 $\lambda'_j$ 是第 $j$ 列高度。当 $j$ 从 $1$ 增加到 $\lambda_1$ 时，这些 hook lengths 构成整数集合
\[
\{1,2,\dots,\ell_1\}
\setminus
\{\ell_1-\ell_2,\ell_1-\ell_3,\dots,\ell_1-\ell_r\}.
\]
确切地说，若跨过某一行末 $j=\lambda_i$，列高会下降一次，于是 hook 序列除通常的减 $1$ 外额外跳过一个整数；第 $i$ 行造成的被跳过值为
\[
\lambda_1+r-1-(\lambda_i+r-i)=\ell_1-\ell_i.
\]
所以
\[
\prod_{j=1}^{\lambda_1}h_{1j}
=\frac{\ell_1!}{\prod_{i=2}^{r}(\ell_1-\ell_i)}.
\]
与归纳假设相乘便得到 \eqnrefs{eq:hook-product-shifted}。
\end{proof}

\begin{lemma}[钩长乘积的角格递推]
若 $\lambda\vdash n$，则
\[
\sum_{\mu\nearrow\lambda}\frac{H_\lambda}{H_\mu}=n.
\]
\end{lemma}

\begin{proof}
先比较删去一个角格前后的钩长乘积。允许在分拆末尾补零，使行数 $r$ 在删格前后保持相同。若从第 $k$ 行末删去角格，则
\[
\ell_k\longmapsto\ell_k-1,
\]
其余 $\ell_i$ 不变。由 \eqnrefs{eq:hook-product-shifted} 直接相除，得到
\begin{align}
\frac{H_\lambda}{H_{\lambda\setminus k}}
&=\ell_k
\prod_{i<k}
\frac{\ell_i-\ell_k+1}{\ell_i-\ell_k}
\prod_{j>k}
\frac{\ell_k-\ell_j-1}{\ell_k-\ell_j}.
\label{eq:hook-corner-ratio}
\end{align}
如果第 $k$ 行不是可移去角，即 $\lambda_k=\lambda_{k+1}$，那么 $\ell_k-\ell_{k+1}=1$，右边第二个乘积自动出现零因子。因此可以把 \eqnrefs{eq:hook-corner-ratio} 对所有 $k=1,\dots,r$ 求和，而非可移去行自动贡献零。

记
\[
A_k=
\prod_{i<k}
\frac{\ell_i-\ell_k+1}{\ell_i-\ell_k}
\prod_{j>k}
\frac{\ell_k-\ell_j-1}{\ell_k-\ell_j}.
\]
则需要证明的关键恒等式是
\begin{equation}
\sum_{k=1}^{r}\ell_kA_k=n.
\label{eq:hook-ratio-sum}
\end{equation}
为此定义有理函数
\[
F(z)=\prod_{i=1}^{r}\frac{z-\ell_i-1}{z-\ell_i}.
\]
它在 $z=\ell_k$ 处只有简单极点，并且
\[
\operatorname*{Res}_{z=\ell_k}F(z)=-A_k.
\]
因此
\begin{align*}
\sum_{k=1}^{r}\ell_kA_k
&=-\sum_{k=1}^{r}\operatorname*{Res}_{z=\ell_k}\bigl[zF(z)\bigr]\\
&=\operatorname*{Res}_{z=\infty}\bigl[zF(z)\bigr],
\end{align*}
其中最后一步使用 Riemann 球面上全部留数之和为零。

只需计算无穷远处的 Laurent 展开。对每个因子，
\[
\frac{z-\ell_i-1}{z-\ell_i}
=1-\frac{1}{z-\ell_i}
=1-\frac1z-\frac{\ell_i}{z^2}+O(z^{-3}).
\]
相乘得
\[
F(z)
=1-\frac rz
+\frac{\binom r2-\sum_i\ell_i}{z^2}
+O(z^{-3}),
\]
所以
\[
zF(z)
=z-r+\frac{\binom r2-\sum_i\ell_i}{z}+O(z^{-2}).
\]
无穷远留数等于 $1/z$ 系数的负值，因此
\[
\operatorname*{Res}_{z=\infty}[zF(z)]
=\sum_i\ell_i-\binom r2.
\]
由 \eqnrefs{eq:shifted-row-lengths}，
\[
\sum_i\ell_i
=\sum_i\lambda_i+\sum_{i=1}^{r}(r-i)
=n+\binom r2,
\]
故无穷远留数正好等于 $n$，从而 \eqnrefs{eq:hook-ratio-sum} 得证。结合 \eqnrefs{eq:hook-corner-ratio}，我们得到
\begin{equation}
\sum_{\mu\nearrow\lambda}\frac{H_\lambda}{H_\mu}=n.
\label{eq:hook-product-recursion}
\end{equation}
\end{proof}

\begin{theorem}[钩长公式与 Young 归一化常数]
对 $\lambda\vdash n$，
\[
f^\lambda=\frac{n!}{H_\lambda},\qquad c_\lambda=H_\lambda,
\qquad p_T=\frac{e_T}{H_\lambda}.
\]
\end{theorem}

\begin{proof}
现在对 $n$ 作归纳。假设所有 $n-1$ 个方格的 Young 图已经满足
\[
f^\mu=\frac{(n-1)!}{H_\mu}.
\]
代入标准盘分支递推 \eqnrefs{eq:hook-recursion-app}，有
\begin{align*}
f^\lambda
&=\sum_{\mu\nearrow\lambda}f^\mu
=(n-1)!\sum_{\mu\nearrow\lambda}\frac1{H_\mu}\\
&=\frac{(n-1)!}{H_\lambda}
\sum_{\mu\nearrow\lambda}\frac{H_\lambda}{H_\mu}\\
&=\frac{n!}{H_\lambda}.
\end{align*}
初始情形 $[1]$ 满足 $f^{[1]}=1=1!/H_{[1]}$，因此归纳闭合：
\[
f^\lambda=\frac{n!}{\displaystyle\prod_{u\in\lambda}h(u)}.
\]

再由 $c_\lambda=n!/\dim Ae_T$ 及 $\dim Ae_T=f^\lambda$，
得到 $c_\lambda=H_\lambda$，归一化常数也随之确定。
\end{proof}

这套证明把 钩长公式拆成三个可以分别核验的步骤：标准盘的分支递推来自“最大数字只能位于可移去角”；钩长乘积的局部变化由 \eqnrefs{eq:hook-corner-ratio} 精确控制；最后的有理函数留数恒等式把所有可移去角的贡献加总为 $n$。因此公式并不是组合计数的猜测，而是一个完整闭合的归纳证明。

\section{从本原幂等元到物理对称化算符}

群代数证明看起来很抽象，但它最终对应一个非常具体的物理操作。若 $S_n$ 作用在多体张量空间 $\mathcal H_1^{\otimes n}$，则把每个置换 $\pi$ 换成粒子标签置换算符 $U(\pi)$，群代数元素
\[
p_T=\frac1{h_\lambda}e_T
\]
就变成一个实际线性算符。由于约定 $e_T=a_Tb_T$，算符从右向左作用：它先在每一列对应的粒子标签间反对称化，再在每一行之间对称化。
$U(p_T)$ 为幂等算符，其像空间具有指定交换对称类型 $\lambda$；
一般它不是正交投影，也不等于整个 $\lambda$ 等型分量的投影。
若需要后者，应使用中心幂等元
\[
P_\lambda=\frac{f^\lambda}{n!}\sum_{\pi\in S_n}
\chi^\lambda(\pi^{-1})U(\pi).
\]

完全对称图 $[n]$ 给出玻色对称化算符
\[
P_+=\frac1{n!}\sum_{\pi\in S_n}U(\pi),
\]
完全反对称图 $[1^n]$ 给出费米反对称化算符
\[
P_-=\frac1{n!}\sum_{\pi\in S_n}\operatorname{sgn}(\pi)U(\pi).
\]
一般 Young 图则插在这两种极端之间，描述多自由度体系中常见的混合交换对称性。于是 Young 方法并不是一套与物理脱节的组合技巧，而是“从群代数本原幂等元构造多体 Hilbert 空间对称扇区”的系统方法。

Yang--Mills 拓扑荷与大统一/seesaw 的物理应用在\chapref{chap:lie-standard-model}继续讨论；下一章先利用这里的标准基与维数公式建立特征标的统一计算方法。
